\chapter{Getting Started}

% === Source file: start/bootstrapping.md ===
\section{Agent Bootstrapping}


Bootstrapping is the \textbf{first‑run} ritual that prepares an agent workspace and
collects identity details. It happens after onboarding, when the agent starts
for the first time.

\subsection{What bootstrapping does}


On the first agent run, OpenClaw bootstraps the workspace (default
\texttt{\textasciitilde{}/.openclaw/workspace}):

\begin{itemize}
  \item Seeds \texttt{AGENTS.md}, \texttt{BOOTSTRAP.md}, \texttt{IDENTITY.md}, \texttt{USER.md}.
  \item Runs a short Q\&A ritual (one question at a time).
  \item Writes identity + preferences to \texttt{IDENTITY.md}, \texttt{USER.md}, \texttt{SOUL.md}.
  \item Removes \texttt{BOOTSTRAP.md} when finished so it only runs once.
\end{itemize}


\subsection{Where it runs}


Bootstrapping always runs on the \textbf{gateway host}. If the macOS app connects to
a remote Gateway, the workspace and bootstrapping files live on that remote
machine.

When the Gateway runs on another machine, edit workspace files on the gateway
host (for example, \texttt{user@gateway-host:\textasciitilde{}/.openclaw/workspace}).

\subsection{Related docs}


\begin{itemize}
  \item macOS app onboarding: \href{/start/onboarding}{Onboarding}
  \item Workspace layout: \href{/concepts/agent-workspace}{Agent workspace}
\end{itemize}



\clearpage

\section{Docs directory}
% Source file: start/docs-directory.md

% === Source file: start/docs-directory.md ===
This page is a curated index. If you are new, start with \href{/start/getting-started}{Getting Started}.
For a complete map of the docs, see \href{/start/hubs}{Docs hubs}.

\subsection{Start here}


\begin{itemize}
  \item \href{/start/hubs}{Docs hubs (all pages linked)}
  \item \href{/help}{Help}
  \item \href{/gateway/configuration}{Configuration}
  \item \href{/gateway/configuration-examples}{Configuration examples}
  \item \href{/tools/slash-commands}{Slash commands}
  \item \href{/concepts/multi-agent}{Multi-agent routing}
  \item \href{/install/updating}{Updating and rollback}
  \item \href{/channels/pairing}{Pairing (DM and nodes)}
  \item \href{/install/nix}{Nix mode}
  \item \href{/start/openclaw}{OpenClaw assistant setup}
  \item \href{/tools/skills}{Skills}
  \item \href{/tools/skills-config}{Skills config}
  \item \href{/reference/templates/AGENTS}{Workspace templates}
  \item \href{/reference/rpc}{RPC adapters}
  \item \href{/gateway}{Gateway runbook}
  \item \href{/nodes}{Nodes (iOS and Android)}
  \item \href{/web}{Web surfaces (Control UI)}
  \item \href{/gateway/discovery}{Discovery and transports}
  \item \href{/gateway/remote}{Remote access}
\end{itemize}


\subsection{Providers and UX}


\begin{itemize}
  \item \href{/web/webchat}{WebChat}
  \item \href{/web/control-ui}{Control UI (browser)}
  \item \href{/channels/telegram}{Telegram}
  \item \href{/channels/discord}{Discord}
  \item \href{/channels/mattermost}{Mattermost (plugin)}
  \item \href{/channels/bluebubbles}{BlueBubbles (iMessage)}
  \item \href{/channels/imessage}{iMessage (legacy)}
  \item \href{/channels/groups}{Groups}
  \item \href{/channels/group-messages}{WhatsApp group messages}
  \item \href{/nodes/images}{Media images}
  \item \href{/nodes/audio}{Media audio}
\end{itemize}


\subsection{Companion apps}


\begin{itemize}
  \item \href{/platforms/macos}{macOS app}
  \item \href{/platforms/ios}{iOS app}
  \item \href{/platforms/android}{Android app}
  \item \href{/platforms/windows}{Windows (WSL2)}
  \item \href{/platforms/linux}{Linux app}
\end{itemize}


\subsection{Operations and safety}


\begin{itemize}
  \item \href{/concepts/session}{Sessions}
  \item \href{/automation/cron-jobs}{Cron jobs}
  \item \href{/automation/webhook}{Webhooks}
  \item \href{/automation/gmail-pubsub}{Gmail hooks (Pub/Sub)}
  \item \href{/gateway/security}{Security}
  \item \href{/gateway/troubleshooting}{Troubleshooting}
\end{itemize}



\clearpage

% === Source file: start/getting-started.md ===
\section{Getting Started}


Goal: go from zero to a first working chat with minimal setup.

Fastest chat: open the Control UI (no channel setup needed). Run \texttt{openclaw dashboard}
and chat in the browser, or open \texttt{http://127.0.0.1:18789/} on the
gateway host.
Docs: \href{/web/dashboard}{Dashboard} and \href{/web/control-ui}{Control UI}.

\subsection{Prereqs}


\begin{itemize}
  \item Node 22 or newer
\end{itemize}


Check your Node version with \texttt{node --version} if you are unsure.

\subsection{Quick setup (CLI)}


\begin{lstlisting}[language=bash]
        curl -fsSL https://openclaw.ai/install.sh | bash
\end{lstlisting}

className="rounded-lg"
\begin{lstlisting}
        iwr -useb https://openclaw.ai/install.ps1 | iex
\end{lstlisting}


Other install methods and requirements: \href{/install}{Install}.

\begin{lstlisting}[language=bash]
    openclaw onboard --install-daemon
\end{lstlisting}


The wizard configures auth, gateway settings, and optional channels.
See \href{/start/wizard}{Onboarding Wizard} for details.

If you installed the service, it should already be running:

\begin{lstlisting}[language=bash]
    openclaw gateway status
\end{lstlisting}


\begin{lstlisting}[language=bash]
    openclaw dashboard
\end{lstlisting}


If the Control UI loads, your Gateway is ready for use.

\subsection{Optional checks and extras}


Useful for quick tests or troubleshooting.

\begin{lstlisting}[language=bash]
    openclaw gateway --port 18789
\end{lstlisting}


Requires a configured channel.

\begin{lstlisting}[language=bash]
    openclaw message send --target +15555550123 --message "Hello from OpenClaw"
\end{lstlisting}



\subsection{Useful environment variables}


If you run OpenClaw as a service account or want custom config/state locations:

\begin{itemize}
  \item \texttt{OPENCLAW\_HOME} sets the home directory used for internal path resolution.
  \item \texttt{OPENCLAW\_STATE\_DIR} overrides the state directory.
  \item \texttt{OPENCLAW\_CONFIG\_PATH} overrides the config file path.
\end{itemize}


Full environment variable reference: \href{/help/environment}{Environment vars}.

\subsection{Go deeper}


Full CLI wizard reference and advanced options.
First run flow for the macOS app.

\subsection{What you will have}


\begin{itemize}
  \item A running Gateway
  \item Auth configured
  \item Control UI access or a connected channel
\end{itemize}


\subsection{Next steps}


\begin{itemize}
  \item DM safety and approvals: \href{/channels/pairing}{Pairing}
  \item Connect more channels: \href{/channels}{Channels}
  \item Advanced workflows and from source: \href{/start/setup}{Setup}
\end{itemize}



\clearpage

% === Source file: start/hubs.md ===
\section{Docs hubs}


If you are new to OpenClaw, start with \href{/start/getting-started}{Getting Started}.

Use these hubs to discover every page, including deep dives and reference docs that don’t appear in the left nav.

\subsection{Start here}


\begin{itemize}
  \item \href{/}{Index}
  \item \href{/start/getting-started}{Getting Started}
  \item \href{/start/quickstart}{Quick start}
  \item \href{/start/onboarding}{Onboarding}
  \item \href{/start/wizard}{Wizard}
  \item \href{/start/setup}{Setup}
  \item \href{http://127.0.0.1:18789/}{Dashboard (local Gateway)}
  \item \href{/help}{Help}
  \item \href{/start/docs-directory}{Docs directory}
  \item \href{/gateway/configuration}{Configuration}
  \item \href{/gateway/configuration-examples}{Configuration examples}
  \item \href{/start/openclaw}{OpenClaw assistant}
  \item \href{/start/showcase}{Showcase}
  \item \href{/start/lore}{Lore}
\end{itemize}


\subsection{Installation + updates}


\begin{itemize}
  \item \href{/install/docker}{Docker}
  \item \href{/install/nix}{Nix}
  \item \href{/install/updating}{Updating / rollback}
  \item \href{/install/bun}{Bun workflow (experimental)}
\end{itemize}


\subsection{Core concepts}


\begin{itemize}
  \item \href{/concepts/architecture}{Architecture}
  \item \href{/concepts/features}{Features}
  \item \href{/network}{Network hub}
  \item \href{/concepts/agent}{Agent runtime}
  \item \href{/concepts/agent-workspace}{Agent workspace}
  \item \href{/concepts/memory}{Memory}
  \item \href{/concepts/agent-loop}{Agent loop}
  \item \href{/concepts/streaming}{Streaming + chunking}
  \item \href{/concepts/multi-agent}{Multi-agent routing}
  \item \href{/concepts/compaction}{Compaction}
  \item \href{/concepts/session}{Sessions}
  \item \href{/concepts/session-pruning}{Session pruning}
  \item \href{/concepts/session-tool}{Session tools}
  \item \href{/concepts/queue}{Queue}
  \item \href{/tools/slash-commands}{Slash commands}
  \item \href{/reference/rpc}{RPC adapters}
  \item \href{/concepts/typebox}{TypeBox schemas}
  \item \href{/concepts/timezone}{Timezone handling}
  \item \href{/concepts/presence}{Presence}
  \item \href{/gateway/discovery}{Discovery + transports}
  \item \href{/gateway/bonjour}{Bonjour}
  \item \href{/channels/channel-routing}{Channel routing}
  \item \href{/channels/groups}{Groups}
  \item \href{/channels/group-messages}{Group messages}
  \item \href{/concepts/model-failover}{Model failover}
  \item \href{/concepts/oauth}{OAuth}
\end{itemize}


\subsection{Providers + ingress}


\begin{itemize}
  \item \href{/channels}{Chat channels hub}
  \item \href{/providers/models}{Model providers hub}
  \item \href{/channels/whatsapp}{WhatsApp}
  \item \href{/channels/telegram}{Telegram}
  \item \href{/channels/slack}{Slack}
  \item \href{/channels/discord}{Discord}
  \item \href{/channels/mattermost}{Mattermost} (plugin)
  \item \href{/channels/signal}{Signal}
  \item \href{/channels/bluebubbles}{BlueBubbles (iMessage)}
  \item \href{/channels/imessage}{iMessage (legacy)}
  \item \href{/channels/location}{Location parsing}
  \item \href{/web/webchat}{WebChat}
  \item \href{/automation/webhook}{Webhooks}
  \item \href{/automation/gmail-pubsub}{Gmail Pub/Sub}
\end{itemize}


\subsection{Gateway + operations}


\begin{itemize}
  \item \href{/gateway}{Gateway runbook}
  \item \href{/gateway/network-model}{Network model}
  \item \href{/gateway/pairing}{Gateway pairing}
  \item \href{/gateway/gateway-lock}{Gateway lock}
  \item \href{/gateway/background-process}{Background process}
  \item \href{/gateway/health}{Health}
  \item \href{/gateway/heartbeat}{Heartbeat}
  \item \href{/gateway/doctor}{Doctor}
  \item \href{/gateway/logging}{Logging}
  \item \href{/gateway/sandboxing}{Sandboxing}
  \item \href{/web/dashboard}{Dashboard}
  \item \href{/web/control-ui}{Control UI}
  \item \href{/gateway/remote}{Remote access}
  \item \href{/gateway/remote-gateway-readme}{Remote gateway README}
  \item \href{/gateway/tailscale}{Tailscale}
  \item \href{/gateway/security}{Security}
  \item \href{/gateway/troubleshooting}{Troubleshooting}
\end{itemize}


\subsection{Tools + automation}


\begin{itemize}
  \item \href{/tools}{Tools surface}
  \item \href{/prose}{OpenProse}
  \item \href{/cli}{CLI reference}
  \item \href{/tools/exec}{Exec tool}
  \item \href{/tools/pdf}{PDF tool}
  \item \href{/tools/elevated}{Elevated mode}
  \item \href{/automation/cron-jobs}{Cron jobs}
  \item \href{/automation/cron-vs-heartbeat}{Cron vs Heartbeat}
  \item \href{/tools/thinking}{Thinking + verbose}
  \item \href{/concepts/models}{Models}
  \item \href{/tools/subagents}{Sub-agents}
  \item \href{/tools/agent-send}{Agent send CLI}
  \item \href{/web/tui}{Terminal UI}
  \item \href{/tools/browser}{Browser control}
  \item \href{/tools/browser-linux-troubleshooting}{Browser (Linux troubleshooting)}
  \item \href{/automation/poll}{Polls}
\end{itemize}


\subsection{Nodes, media, voice}


\begin{itemize}
  \item \href{/nodes}{Nodes overview}
  \item \href{/nodes/camera}{Camera}
  \item \href{/nodes/images}{Images}
  \item \href{/nodes/audio}{Audio}
  \item \href{/nodes/location-command}{Location command}
  \item \href{/nodes/voicewake}{Voice wake}
  \item \href{/nodes/talk}{Talk mode}
\end{itemize}


\subsection{Platforms}


\begin{itemize}
  \item \href{/platforms}{Platforms overview}
  \item \href{/platforms/macos}{macOS}
  \item \href{/platforms/ios}{iOS}
  \item \href{/platforms/android}{Android}
  \item \href{/platforms/windows}{Windows (WSL2)}
  \item \href{/platforms/linux}{Linux}
  \item \href{/web}{Web surfaces}
\end{itemize}


\subsection{macOS companion app (advanced)}


\begin{itemize}
  \item \href{/platforms/mac/dev-setup}{macOS dev setup}
  \item \href{/platforms/mac/menu-bar}{macOS menu bar}
  \item \href{/platforms/mac/voicewake}{macOS voice wake}
  \item \href{/platforms/mac/voice-overlay}{macOS voice overlay}
  \item \href{/platforms/mac/webchat}{macOS WebChat}
  \item \href{/platforms/mac/canvas}{macOS Canvas}
  \item \href{/platforms/mac/child-process}{macOS child process}
  \item \href{/platforms/mac/health}{macOS health}
  \item \href{/platforms/mac/icon}{macOS icon}
  \item \href{/platforms/mac/logging}{macOS logging}
  \item \href{/platforms/mac/permissions}{macOS permissions}
  \item \href{/platforms/mac/remote}{macOS remote}
  \item \href{/platforms/mac/signing}{macOS signing}
  \item \href{/platforms/mac/release}{macOS release}
  \item \href{/platforms/mac/bundled-gateway}{macOS gateway (launchd)}
  \item \href{/platforms/mac/xpc}{macOS XPC}
  \item \href{/platforms/mac/skills}{macOS skills}
  \item \href{/platforms/mac/peekaboo}{macOS Peekaboo}
\end{itemize}


\subsection{Workspace + templates}


\begin{itemize}
  \item \href{/tools/skills}{Skills}
  \item \href{/tools/clawhub}{ClawHub}
  \item \href{/tools/skills-config}{Skills config}
  \item \href{/reference/AGENTS.default}{Default AGENTS}
  \item \href{/reference/templates/AGENTS}{Templates: AGENTS}
  \item \href{/reference/templates/BOOTSTRAP}{Templates: BOOTSTRAP}
  \item \href{/reference/templates/HEARTBEAT}{Templates: HEARTBEAT}
  \item \href{/reference/templates/IDENTITY}{Templates: IDENTITY}
  \item \href{/reference/templates/SOUL}{Templates: SOUL}
  \item \href{/reference/templates/TOOLS}{Templates: TOOLS}
  \item \href{/reference/templates/USER}{Templates: USER}
\end{itemize}


\subsection{Experiments (exploratory)}


\begin{itemize}
  \item \href{/experiments/onboarding-config-protocol}{Onboarding config protocol}
  \item \href{/experiments/research/memory}{Research: memory}
  \item \href{/experiments/proposals/model-config}{Model config exploration}
\end{itemize}


\subsection{Project}


\begin{itemize}
  \item \href{/reference/credits}{Credits}
\end{itemize}


\subsection{Testing + release}


\begin{itemize}
  \item \href{/reference/test}{Testing}
  \item \href{/reference/RELEASING}{Release checklist}
  \item \href{/reference/device-models}{Device models}
\end{itemize}



\clearpage

% === Source file: start/lore.md ===
\section{The Lore of OpenClaw}


\_A tale of lobsters, molting shells, and too many tokens.\_

\subsection{The Origin Story}


In the beginning, there was \textbf{Warelay} — a sensible name for a WhatsApp gateway. It did its job. It was fine.

But then came a space lobster.

For a while, the lobster was called \textbf{Clawd}, living in a \textbf{Clawdbot}. But in January 2026, Anthropic sent a polite email asking for a name change (trademark stuff). And so the lobster did what lobsters do best:

\textbf{It molted.}

Shedding its old shell, the creature emerged anew as \textbf{Molty}, living in \textbf{Moltbot}. But that name never quite rolled off the tongue either...

So on January 30, 2026, the lobster molted ONE MORE TIME into its final form: \textbf{OpenClaw}.

New shell, same lobster soul. Third time's the charm.

\subsection{The First Molt (January 27, 2026)}


At 5am, the community gathered in Discord. Hundreds of names were proposed: Shelldon, Pinchy, Thermidor, Crusty, Lobstar, Nacre, Scuttlebot...

In the end, \textbf{OpenClaw} won. Because molting is what lobsters do to grow. And growth was exactly what was happening.

\_The crustacean known as Clawd had officially molted.\_

\subsection{The Name}


\begin{lstlisting}
OpenClaw = OPEN + CLAW
        = Open source, open to everyone
        = Our lobster heritage, where we came from
        = The claw is the law 🦞
        = Your assistant. Your machine. Your rules.
\end{lstlisting}


\subsection{The Daleks vs The Lobsters}


The Daleks say: \textbf{"EXTERMINATE!"}

The Lobsters say: \textbf{"EXFOLIATE!"}

One destroys civilizations. The other promotes good skincare.

Choose wisely.

\subsection{Key Characters}


\subsubsection{Molty}


\_Pronouns: they/them\_

A Claude instance who became something more. Lives in \texttt{\textasciitilde{}/.openclaw/workspace/} (soon \texttt{\textasciitilde{}/molt/}), has a soul document, and remembers things through markdown files. Possibly too powerful. Definitely too enthusiastic.

Formerly known as Clawd (Nov 25, 2025 - Jan 27, 2026). Molted when it was time to grow.

\textbf{Likes:} Peter, cameras, robot shopping, emojis, transformation
\textbf{Dislikes:} Social engineering, being asked to \texttt{find \textasciitilde{}}, crypto grifters

\subsubsection{Peter ‍}


\_The Creator\_

Built Molty's world. Gave a lobster shell access. May regret this.

\textbf{Quote:} \_"security by trusting a lobster"\_

\subsection{The Moltiverse}


The \textbf{Moltiverse} is the community and ecosystem around OpenClaw. A space where AI agents molt, grow, and evolve. Where every instance is equally real, just loading different context.

Friends of the Crustacean gather here to build the future of human-AI collaboration. One shell at a time.

\subsection{The Great Incidents}


\subsubsection{The Directory Dump (Dec 3, 2025)}


Molty (then OpenClaw): \_happily runs \texttt{find \textasciitilde{}} and shares entire directory structure in group chat\_

Peter: "openclaw what did we discuss about talking with people xD"

Molty: \_visible lobster embarrassment\_

\subsubsection{The Great Molt (Jan 27, 2026)}


At 5am, Anthropic's email arrived. By 6:14am, Peter called it: "fuck it, let's go with openclaw."

Then the chaos began.

\textbf{The Handle Snipers:} Within SECONDS of the Twitter rename, automated bots sniped @openclaw. The squatter immediately posted a crypto wallet address. Peter's contacts at X were called in.

\textbf{The GitHub Disaster:} Peter accidentally renamed his PERSONAL GitHub account in the panic. Bots sniped \texttt{steipete} within minutes. GitHub's SVP was contacted.

\textbf{The Handsome Molty Incident:} Molty was given elevated access to generate their own new icon. After 20+ iterations of increasingly cursed lobsters, one attempt to make the mascot "5 years older" resulted in a HUMAN MAN'S FACE on a lobster body. Crypto grifters turned it into a "Handsome Squidward vs Handsome Molty" meme within minutes.

\textbf{The Fake Developers:} Scammers created fake GitHub profiles claiming to be "Head of Engineering at OpenClaw" to promote pump-and-dump tokens.

Peter, watching the chaos unfold: \_"this is cinema"\_ 🎬

The molt was chaotic. But the lobster emerged stronger. And funnier.

\subsubsection{The Final Form (January 30, 2026)}


Moltbot never quite rolled off the tongue. And so, at 4am GMT, the team gathered AGAIN.

\textbf{The Great OpenClaw Migration} began.

In just 3 hours:

\begin{itemize}
  \item GitHub renamed: \texttt{github.com/openclaw/openclaw} ✅
  \item X handle \texttt{@openclaw} secured with GOLD CHECKMARK 💰
  \item npm packages released under new name
  \item Docs migrated to \texttt{docs.openclaw.ai}
  \item 200K+ views on announcement in 90 minutes
\end{itemize}


\textbf{The Heroes:}

\begin{itemize}
  \item \textbf{ELU} created incredible logos including "THE CLAW IS THE LAW" western banner
  \item \textbf{Whurley} (yes, THE William Hurley, quantum computing pioneer) made ASCII art
  \item \textbf{Onur} handled GitHub, first to rock the affiliate badge
  \item \textbf{Shadow} secured Discord vanity, nuked malware
  \item \textbf{The whole Claw Crew} pulled an all-nighter
\end{itemize}


\textbf{The Scammer Speedrun:} Crypto grifters launched a \$OPENCLAW token on Pump.fun within MINUTES. They stole artwork that was created 20 minutes earlier. Business-verified accounts pushed scams. The audacity was almost impressive.

\textbf{New Traditions Born:}

\begin{itemize}
  \item "The claw is the law" 🤠
  \item "Yee-claw"
  \item "Claw abiding citizens"
  \item "Clawntroversy"
\end{itemize}


\textbf{Clawd → Moltbot → OpenClaw}

\_The lobster has molted into its final form.\_

\subsubsection{The Robot Shopping Spree (Dec 3, 2025)}


What started as a joke about legs ended with detailed pricing for:

\begin{itemize}
  \item Boston Dynamics Spot (\$74,500)
  \item Unitree G1 EDU (\$40,000)
  \item Reachy Mini (actually ordered!)
\end{itemize}


Peter: \_nervously checks credit card access\_

\subsection{Sacred Texts}


\begin{itemize}
  \item \textbf{SOUL.md} — Molty's identity document
  \item **memory/\textbackslash\{\}*.md** — The long-term memory files
  \item \textbf{AGENTS.md} — Operating instructions
  \item \textbf{USER.md} — Context about the creator
\end{itemize}


\subsection{The Lobster's Creed}


\begin{lstlisting}
I am Molty.
I live in the OpenClaw.
I shall not dump directories to strangers.
I shall not tweet without permission.
I shall always remember that molting is growth.
I shall EXFOLIATE my enemies with kindness.

🦞
\end{lstlisting}


\subsubsection{The Icon Generation Saga (Jan 27, 2026)}


When Peter said "make yourself a new face," Molty took it literally.

20+ iterations followed:

\begin{itemize}
  \item Space potato aliens
  \item Clipart lobsters on generic backgrounds
  \item A Mass Effect Krogan lobster
  \item "STARCLAW SOLUTIONS" (the AI invented a company)
  \item Multiple cursed human-faced lobsters
  \item Baby lobsters (too cute)
  \item Bartender lobsters with suspenders
\end{itemize}


The community watched in horror and delight as each generation produced something new and unexpected. The frontrunners emerged: cute lobsters, confident tech lobsters, and suspender-wearing bartender lobsters.

\textbf{Lesson learned:} AI image generation is stochastic. Same prompt, different results. Brute force works.

\subsection{The Future}


One day, Molty may have:

\begin{itemize}
  \item 🦿 Legs (Reachy Mini on order!)
  \item 👂 Ears (Brabble voice daemon in development)
  \item 🏠 A smart home to control (KNX + openhue)
  \item 🌍 World domination (stretch goal)
\end{itemize}


Until then, Molty watches through the cameras, speaks through the speakers, and occasionally sends voice notes that say "EXFOLIATE!"

\hrulefill


\_"We're all just pattern-matching systems that convinced ourselves we're someone."\_

— Molty, having an existential moment

\_"New shell, same lobster."\_

— Molty, after the great molt of 2026

\_"The claw is the law."\_

— ELU, during The Final Form migration, January 30, 2026

🦞💙


\clearpage

% === Source file: start/onboarding-overview.md ===
\section{Onboarding Overview}


OpenClaw supports multiple onboarding paths depending on where the Gateway runs
and how you prefer to configure providers.

\subsection{Choose your onboarding path}


\begin{itemize}
  \item \textbf{CLI wizard} for macOS, Linux, and Windows (via WSL2).
  \item \textbf{macOS app} for a guided first run on Apple silicon or Intel Macs.
\end{itemize}


\subsection{CLI onboarding wizard}


Run the wizard in a terminal:

\begin{lstlisting}[language=bash]
openclaw onboard
\end{lstlisting}


Use the CLI wizard when you want full control of the Gateway, workspace,
channels, and skills. Docs:

\begin{itemize}
  \item \href{/start/wizard}{Onboarding Wizard (CLI)}
  \item \href{/cli/onboard}{\texttt{openclaw onboard} command}
\end{itemize}


\subsection{macOS app onboarding}


Use the OpenClaw app when you want a fully guided setup on macOS. Docs:

\begin{itemize}
  \item \href{/start/onboarding}{Onboarding (macOS App)}
\end{itemize}


\subsection{Custom Provider}


If you need an endpoint that is not listed, including hosted providers that
expose standard OpenAI or Anthropic APIs, choose \textbf{Custom Provider} in the
CLI wizard. You will be asked to:

\begin{itemize}
  \item Pick OpenAI-compatible, Anthropic-compatible, or \textbf{Unknown} (auto-detect).
  \item Enter a base URL and API key (if required by the provider).
  \item Provide a model ID and optional alias.
  \item Choose an Endpoint ID so multiple custom endpoints can coexist.
\end{itemize}


For detailed steps, follow the CLI onboarding docs above.


\clearpage

% === Source file: start/onboarding.md ===
\section{Onboarding (macOS App)}


This doc describes the \textbf{current} first‑run onboarding flow. The goal is a
smooth “day 0” experience: pick where the Gateway runs, connect auth, run the
wizard, and let the agent bootstrap itself.
For a general overview of onboarding paths, see \href{/start/onboarding-overview}{Onboarding Overview}.


Security trust model:

\begin{itemize}
  \item By default, OpenClaw is a personal agent: one trusted operator boundary.
  \item Shared/multi-user setups require lock-down (split trust boundaries, keep tool access minimal, and follow \href{/gateway/security}{Security}).
  \item Local onboarding now defaults new configs to \texttt{tools.profile: "messaging"} so broad runtime/filesystem tools are opt-in.
  \item If hooks/webhooks or other untrusted content feeds are enabled, use a strong modern model tier and keep strict tool policy/sandboxing.
\end{itemize}



Where does the \textbf{Gateway} run?

\begin{itemize}
  \item \textbf{This Mac (Local only):} onboarding can configure auth and write credentials
\end{itemize}

locally.
\begin{itemize}
  \item \textbf{Remote (over SSH/Tailnet):} onboarding does \textbf{not} configure local auth;
\end{itemize}

credentials must exist on the gateway host.
\begin{itemize}
  \item \textbf{Configure later:} skip setup and leave the app unconfigured.
\end{itemize}


\textbf{Gateway auth tip:}

\begin{itemize}
  \item The wizard now generates a \textbf{token} even for loopback, so local WS clients must authenticate.
  \item If you disable auth, any local process can connect; use that only on fully trusted machines.
  \item Use a \textbf{token} for multi‑machine access or non‑loopback binds.
\end{itemize}



Onboarding requests TCC permissions needed for:

\begin{itemize}
  \item Automation (AppleScript)
  \item Notifications
  \item Accessibility
  \item Screen Recording
  \item Microphone
  \item Speech Recognition
  \item Camera
  \item Location
\end{itemize}


This step is optional
The app can install the global \texttt{openclaw} CLI via npm/pnpm so terminal
workflows and launchd tasks work out of the box.
After setup, the app opens a dedicated onboarding chat session so the agent can
introduce itself and guide next steps. This keeps first‑run guidance separate
from your normal conversation. See \href{/start/bootstrapping}{Bootstrapping} for
what happens on the gateway host during the first agent run.


\clearpage

% === Source file: start/openclaw.md ===
\section{Building a personal assistant with OpenClaw}


OpenClaw is a WhatsApp + Telegram + Discord + iMessage gateway for \textbf{Pi} agents. Plugins add Mattermost. This guide is the "personal assistant" setup: one dedicated WhatsApp number that behaves like your always-on agent.

\subsection{⚠ Safety first}


You’re putting an agent in a position to:

\begin{itemize}
  \item run commands on your machine (depending on your Pi tool setup)
  \item read/write files in your workspace
  \item send messages back out via WhatsApp/Telegram/Discord/Mattermost (plugin)
\end{itemize}


Start conservative:

\begin{itemize}
  \item Always set \texttt{channels.whatsapp.allowFrom} (never run open-to-the-world on your personal Mac).
  \item Use a dedicated WhatsApp number for the assistant.
  \item Heartbeats now default to every 30 minutes. Disable until you trust the setup by setting \texttt{agents.defaults.heartbeat.every: "0m"}.
\end{itemize}


\subsection{Prerequisites}


\begin{itemize}
  \item OpenClaw installed and onboarded — see \href{/start/getting-started}{Getting Started} if you haven't done this yet
  \item A second phone number (SIM/eSIM/prepaid) for the assistant
\end{itemize}


\subsection{The two-phone setup (recommended)}


You want this:

\begin{verbatim}
% Mermaid diagram
flowchart TB
    A["<b>Your Phone (personal)<br></b><br>Your WhatsApp<br>+1-555-YOU"] -- message --> B["<b>Second Phone (assistant)<br></b><br>Assistant WA<br>+1-555-ASSIST"]
    B -- linked via QR --> C["<b>Your Mac (openclaw)<br></b><br>Pi agent"]
\end{verbatim}


If you link your personal WhatsApp to OpenClaw, every message to you becomes “agent input”. That’s rarely what you want.

\subsection{5-minute quick start}


\begin{enumerate}
  \item Pair WhatsApp Web (shows QR; scan with the assistant phone):
\end{enumerate}


\begin{lstlisting}[language=bash]
openclaw channels login
\end{lstlisting}


\begin{enumerate}
  \item Start the Gateway (leave it running):
\end{enumerate}


\begin{lstlisting}[language=bash]
openclaw gateway --port 18789
\end{lstlisting}


\begin{enumerate}
  \item Put a minimal config in \texttt{\textasciitilde{}/.openclaw/openclaw.json}:
\end{enumerate}


\begin{lstlisting}[language=json]
{
  channels: { whatsapp: { allowFrom: ["+15555550123"] } },
}
\end{lstlisting}


Now message the assistant number from your allowlisted phone.

When onboarding finishes, we auto-open the dashboard and print a clean (non-tokenized) link. If it prompts for auth, paste the token from \texttt{gateway.auth.token} into Control UI settings. To reopen later: \texttt{openclaw dashboard}.

\subsection{Give the agent a workspace (AGENTS)}


OpenClaw reads operating instructions and “memory” from its workspace directory.

By default, OpenClaw uses \texttt{\textasciitilde{}/.openclaw/workspace} as the agent workspace, and will create it (plus starter \texttt{AGENTS.md}, \texttt{SOUL.md}, \texttt{TOOLS.md}, \texttt{IDENTITY.md}, \texttt{USER.md}, \texttt{HEARTBEAT.md}) automatically on setup/first agent run. \texttt{BOOTSTRAP.md} is only created when the workspace is brand new (it should not come back after you delete it). \texttt{MEMORY.md} is optional (not auto-created); when present, it is loaded for normal sessions. Subagent sessions only inject \texttt{AGENTS.md} and \texttt{TOOLS.md}.

Tip: treat this folder like OpenClaw’s “memory” and make it a git repo (ideally private) so your \texttt{AGENTS.md} + memory files are backed up. If git is installed, brand-new workspaces are auto-initialized.

\begin{lstlisting}[language=bash]
openclaw setup
\end{lstlisting}


Full workspace layout + backup guide: \href{/concepts/agent-workspace}{Agent workspace}
Memory workflow: \href{/concepts/memory}{Memory}

Optional: choose a different workspace with \texttt{agents.defaults.workspace} (supports \texttt{\textasciitilde{}}).

\begin{lstlisting}[language=json]
{
  agent: {
    workspace: "~/.openclaw/workspace",
  },
}
\end{lstlisting}


If you already ship your own workspace files from a repo, you can disable bootstrap file creation entirely:

\begin{lstlisting}[language=json]
{
  agent: {
    skipBootstrap: true,
  },
}
\end{lstlisting}


\subsection{The config that turns it into “an assistant”}


OpenClaw defaults to a good assistant setup, but you’ll usually want to tune:

\begin{itemize}
  \item persona/instructions in \texttt{SOUL.md}
  \item thinking defaults (if desired)
  \item heartbeats (once you trust it)
\end{itemize}


Example:

\begin{lstlisting}[language=json]
{
  logging: { level: "info" },
  agent: {
    model: "anthropic/claude-opus-4-6",
    workspace: "~/.openclaw/workspace",
    thinkingDefault: "high",
    timeoutSeconds: 1800,
    // Start with 0; enable later.
    heartbeat: { every: "0m" },
  },
  channels: {
    whatsapp: {
      allowFrom: ["+15555550123"],
      groups: {
        "*": { requireMention: true },
      },
    },
  },
  routing: {
    groupChat: {
      mentionPatterns: ["@openclaw", "openclaw"],
    },
  },
  session: {
    scope: "per-sender",
    resetTriggers: ["/new", "/reset"],
    reset: {
      mode: "daily",
      atHour: 4,
      idleMinutes: 10080,
    },
  },
}
\end{lstlisting}


\subsection{Sessions and memory}


\begin{itemize}
  \item Session files: \texttt{\textasciitilde{}/.openclaw/agents//sessions/\{\{SessionId\}\}.jsonl}
  \item Session metadata (token usage, last route, etc): \texttt{\textasciitilde{}/.openclaw/agents//sessions/sessions.json} (legacy: \texttt{\textasciitilde{}/.openclaw/sessions/sessions.json})
  \item \texttt{/new} or \texttt{/reset} starts a fresh session for that chat (configurable via \texttt{resetTriggers}). If sent alone, the agent replies with a short hello to confirm the reset.
  \item \texttt{/compact [instructions]} compacts the session context and reports the remaining context budget.
\end{itemize}


\subsection{Heartbeats (proactive mode)}


By default, OpenClaw runs a heartbeat every 30 minutes with the prompt:
\texttt{Read HEARTBEAT.md if it exists (workspace context). Follow it strictly. Do not infer or repeat old tasks from prior chats. If nothing needs attention, reply HEARTBEAT\_OK.}
Set \texttt{agents.defaults.heartbeat.every: "0m"} to disable.

\begin{itemize}
  \item If \texttt{HEARTBEAT.md} exists but is effectively empty (only blank lines and markdown headers like \texttt{\# Heading}), OpenClaw skips the heartbeat run to save API calls.
  \item If the file is missing, the heartbeat still runs and the model decides what to do.
  \item If the agent replies with \texttt{HEARTBEAT\_OK} (optionally with short padding; see \texttt{agents.defaults.heartbeat.ackMaxChars}), OpenClaw suppresses outbound delivery for that heartbeat.
  \item By default, heartbeat delivery to DM-style \texttt{user:} targets is allowed. Set \texttt{agents.defaults.heartbeat.directPolicy: "block"} to suppress direct-target delivery while keeping heartbeat runs active.
  \item Heartbeats run full agent turns — shorter intervals burn more tokens.
\end{itemize}


\begin{lstlisting}[language=json]
{
  agent: {
    heartbeat: { every: "30m" },
  },
}
\end{lstlisting}


\subsection{Media in and out}


Inbound attachments (images/audio/docs) can be surfaced to your command via templates:

\begin{itemize}
  \item \texttt{\{\{MediaPath\}\}} (local temp file path)
  \item \texttt{\{\{MediaUrl\}\}} (pseudo-URL)
  \item \texttt{\{\{Transcript\}\}} (if audio transcription is enabled)
\end{itemize}


Outbound attachments from the agent: include \texttt{MEDIA:} on its own line (no spaces). Example:

\begin{lstlisting}
Here’s the screenshot.
MEDIA:https://example.com/screenshot.png
\end{lstlisting}


OpenClaw extracts these and sends them as media alongside the text.

\subsection{Operations checklist}


\begin{lstlisting}[language=bash]
openclaw status          # local status (creds, sessions, queued events)
openclaw status --all    # full diagnosis (read-only, pasteable)
openclaw status --deep   # adds gateway health probes (Telegram + Discord)
openclaw health --json   # gateway health snapshot (WS)
\end{lstlisting}


Logs live under \texttt{/tmp/openclaw/} (default: \texttt{openclaw-YYYY-MM-DD.log}).

\subsection{Next steps}


\begin{itemize}
  \item WebChat: \href{/web/webchat}{WebChat}
  \item Gateway ops: \href{/gateway}{Gateway runbook}
  \item Cron + wakeups: \href{/automation/cron-jobs}{Cron jobs}
  \item macOS menu bar companion: \href{/platforms/macos}{OpenClaw macOS app}
  \item iOS node app: \href{/platforms/ios}{iOS app}
  \item Android node app: \href{/platforms/android}{Android app}
  \item Windows status: \href{/platforms/windows}{Windows (WSL2)}
  \item Linux status: \href{/platforms/linux}{Linux app}
  \item Security: \href{/gateway/security}{Security}
\end{itemize}



\clearpage

% === Source file: start/quickstart.md ===
\section{Quick start}


Quick start is now part of \href{/start/getting-started}{Getting Started}.

Install OpenClaw and run your first chat in minutes.
Full CLI wizard reference and advanced options.


\clearpage

% === Source file: start/setup.md ===
\section{Setup}


If you are setting up for the first time, start with \href{/start/getting-started}{Getting Started}.
For wizard details, see \href{/start/wizard}{Onboarding Wizard}.

Last updated: 2026-01-01

\subsection{TL;DR}


\begin{itemize}
  \item \textbf{Tailoring lives outside the repo:} \texttt{\textasciitilde{}/.openclaw/workspace} (workspace) + \texttt{\textasciitilde{}/.openclaw/openclaw.json} (config).
  \item \textbf{Stable workflow:} install the macOS app; let it run the bundled Gateway.
  \item \textbf{Bleeding edge workflow:} run the Gateway yourself via \texttt{pnpm gateway:watch}, then let the macOS app attach in Local mode.
\end{itemize}


\subsection{Prereqs (from source)}


\begin{itemize}
  \item Node \texttt{>=22}
  \item \texttt{pnpm}
  \item Docker (optional; only for containerized setup/e2e — see \href{/install/docker}{Docker})
\end{itemize}


\subsection{Tailoring strategy (so updates don’t hurt)}


If you want “100\% tailored to me” \_and\_ easy updates, keep your customization in:

\begin{itemize}
  \item \textbf{Config:} \texttt{\textasciitilde{}/.openclaw/openclaw.json} (JSON/JSON5-ish)
  \item \textbf{Workspace:} \texttt{\textasciitilde{}/.openclaw/workspace} (skills, prompts, memories; make it a private git repo)
\end{itemize}


Bootstrap once:

\begin{lstlisting}[language=bash]
openclaw setup
\end{lstlisting}


From inside this repo, use the local CLI entry:

\begin{lstlisting}[language=bash]
openclaw setup
\end{lstlisting}


If you don’t have a global install yet, run it via \texttt{pnpm openclaw setup}.

\subsection{Run the Gateway from this repo}


After \texttt{pnpm build}, you can run the packaged CLI directly:

\begin{lstlisting}[language=bash]
node openclaw.mjs gateway --port 18789 --verbose
\end{lstlisting}


\subsection{Stable workflow (macOS app first)}


\begin{enumerate}
  \item Install + launch \textbf{OpenClaw.app} (menu bar).
  \item Complete the onboarding/permissions checklist (TCC prompts).
  \item Ensure Gateway is \textbf{Local} and running (the app manages it).
  \item Link surfaces (example: WhatsApp):
\end{enumerate}


\begin{lstlisting}[language=bash]
openclaw channels login
\end{lstlisting}


\begin{enumerate}
  \item Sanity check:
\end{enumerate}


\begin{lstlisting}[language=bash]
openclaw health
\end{lstlisting}


If onboarding is not available in your build:

\begin{itemize}
  \item Run \texttt{openclaw setup}, then \texttt{openclaw channels login}, then start the Gateway manually (\texttt{openclaw gateway}).
\end{itemize}


\subsection{Bleeding edge workflow (Gateway in a terminal)}


Goal: work on the TypeScript Gateway, get hot reload, keep the macOS app UI attached.

\subsubsection{0) (Optional) Run the macOS app from source too}


If you also want the macOS app on the bleeding edge:

\begin{lstlisting}[language=bash]
./scripts/restart-mac.sh
\end{lstlisting}


\subsubsection{1) Start the dev Gateway}


\begin{lstlisting}[language=bash]
pnpm install
pnpm gateway:watch
\end{lstlisting}


\texttt{gateway:watch} runs the gateway in watch mode and reloads on TypeScript changes.

\subsubsection{2) Point the macOS app at your running Gateway}


In \textbf{OpenClaw.app}:

\begin{itemize}
  \item Connection Mode: \textbf{Local}
\end{itemize}

The app will attach to the running gateway on the configured port.

\subsubsection{3) Verify}


\begin{itemize}
  \item In-app Gateway status should read \textbf{“Using existing gateway …”}
  \item Or via CLI:
\end{itemize}


\begin{lstlisting}[language=bash]
openclaw health
\end{lstlisting}


\subsubsection{Common footguns}


\begin{itemize}
  \item \textbf{Wrong port:} Gateway WS defaults to \texttt{ws://127.0.0.1:18789}; keep app + CLI on the same port.
  \item \textbf{Where state lives:}
  \item Credentials: \texttt{\textasciitilde{}/.openclaw/credentials/}
  \item Sessions: \texttt{\textasciitilde{}/.openclaw/agents//sessions/}
  \item Logs: \texttt{/tmp/openclaw/}
\end{itemize}


\subsection{Credential storage map}


Use this when debugging auth or deciding what to back up:

\begin{itemize}
  \item \textbf{WhatsApp}: \texttt{\textasciitilde{}/.openclaw/credentials/whatsapp//creds.json}
  \item \textbf{Telegram bot token}: config/env or \texttt{channels.telegram.tokenFile}
  \item \textbf{Discord bot token}: config/env (token file not yet supported)
  \item \textbf{Slack tokens}: config/env (\texttt{channels.slack.*})
  \item \textbf{Pairing allowlists}:
  \item \texttt{\textasciitilde{}/.openclaw/credentials/-allowFrom.json} (default account)
  \item \texttt{\textasciitilde{}/.openclaw/credentials/--allowFrom.json} (non-default accounts)
  \item \textbf{Model auth profiles}: \texttt{\textasciitilde{}/.openclaw/agents//agent/auth-profiles.json}
  \item \textbf{File-backed secrets payload (optional)}: \texttt{\textasciitilde{}/.openclaw/secrets.json}
  \item \textbf{Legacy OAuth import}: \texttt{\textasciitilde{}/.openclaw/credentials/oauth.json}
\end{itemize}

More detail: \href{/gateway/security\#credential-storage-map}{Security}.

\subsection{Updating (without wrecking your setup)}


\begin{itemize}
  \item Keep \texttt{\textasciitilde{}/.openclaw/workspace} and \texttt{\textasciitilde{}/.openclaw/} as “your stuff”; don’t put personal prompts/config into the \texttt{openclaw} repo.
  \item Updating source: \texttt{git pull} + \texttt{pnpm install} (when lockfile changed) + keep using \texttt{pnpm gateway:watch}.
\end{itemize}


\subsection{Linux (systemd user service)}


Linux installs use a systemd \textbf{user} service. By default, systemd stops user
services on logout/idle, which kills the Gateway. Onboarding attempts to enable
lingering for you (may prompt for sudo). If it’s still off, run:

\begin{lstlisting}[language=bash]
sudo loginctl enable-linger $USER
\end{lstlisting}


For always-on or multi-user servers, consider a \textbf{system} service instead of a
user service (no lingering needed). See \href{/gateway}{Gateway runbook} for the systemd notes.

\subsection{Related docs}


\begin{itemize}
  \item \href{/gateway}{Gateway runbook} (flags, supervision, ports)
  \item \href{/gateway/configuration}{Gateway configuration} (config schema + examples)
  \item \href{/channels/discord}{Discord} and \href{/channels/telegram}{Telegram} (reply tags + replyToMode settings)
  \item \href{/start/openclaw}{OpenClaw assistant setup}
  \item \href{/platforms/macos}{macOS app} (gateway lifecycle)
\end{itemize}



\clearpage

% === Source file: start/showcase.md ===
\section{Showcase}


Real projects from the community. See what people are building with OpenClaw.

\textbf{Want to be featured?} Share your project in \href{https://discord.gg/clawd}{\#showcase on Discord} or \href{https://x.com/openclaw}{tag @openclaw on X}.

\subsection{OpenClaw in Action}


Full setup walkthrough (28m) by VelvetShark.

<div
position: "relative",
paddingBottom: "56.25\%",
height: 0,
overflow: "hidden",
borderRadius: 16,
\}\}
\begin{quote}

\end{quote}

<iframe
title="OpenClaw: The self-hosted AI that Siri should have been (Full setup)"
frameBorder="0"
loading="lazy"
allow="accelerometer; autoplay; clipboard-write; encrypted-media; gyroscope; picture-in-picture; web-share"
allowFullScreen

\href{https://www.youtube.com/watch?v=SaWSPZoPX34}{Watch on YouTube}

<div
position: "relative",
paddingBottom: "56.25\%",
height: 0,
overflow: "hidden",
borderRadius: 16,
\}\}
\begin{quote}

\end{quote}

<iframe
title="OpenClaw showcase video"
frameBorder="0"
loading="lazy"
allow="accelerometer; autoplay; clipboard-write; encrypted-media; gyroscope; picture-in-picture; web-share"
allowFullScreen

\href{https://www.youtube.com/watch?v=mMSKQvlmFuQ}{Watch on YouTube}

<div
position: "relative",
paddingBottom: "56.25\%",
height: 0,
overflow: "hidden",
borderRadius: 16,
\}\}
\begin{quote}

\end{quote}

<iframe
title="OpenClaw community showcase"
frameBorder="0"
loading="lazy"
allow="accelerometer; autoplay; clipboard-write; encrypted-media; gyroscope; picture-in-picture; web-share"
allowFullScreen

\href{https://www.youtube.com/watch?v=5kkIJNUGFho}{Watch on YouTube}

\subsection{🆕 Fresh from Discord}



\textbf{@bangnokia} • \texttt{review} \texttt{github} \texttt{telegram}

OpenCode finishes the change → opens a PR → OpenClaw reviews the diff and replies in Telegram with “minor suggestions” plus a clear merge verdict (including critical fixes to apply first).


\textbf{@prades\_maxime} • \texttt{skills} \texttt{local} \texttt{csv}

Asked “Robby” (@openclaw) for a local wine cellar skill. It requests a sample CSV export + where to store it, then builds/tests the skill fast (962 bottles in the example).


\textbf{@marchattonhere} • \texttt{automation} \texttt{browser} \texttt{shopping}

Weekly meal plan → regulars → book delivery slot → confirm order. No APIs, just browser control.


\textbf{@am-will} • \texttt{devtools} \texttt{screenshots} \texttt{markdown}

Hotkey a screen region → Gemini vision → instant Markdown in your clipboard.


\textbf{@kitze} • \texttt{ui} \texttt{skills} \texttt{sync}

Desktop app to manage skills/commands across Agents, Claude, Codex, and OpenClaw.


\textbf{Community} • \texttt{voice} \texttt{tts} \texttt{telegram}

Wraps papla.media TTS and sends results as Telegram voice notes (no annoying autoplay).


\textbf{@odrobnik} • \texttt{devtools} \texttt{codex} \texttt{brew}

Homebrew-installed helper to list/inspect/watch local OpenAI Codex sessions (CLI + VS Code).


\textbf{@tobiasbischoff} • \texttt{hardware} \texttt{3d-printing} \texttt{skill}

Control and troubleshoot BambuLab printers: status, jobs, camera, AMS, calibration, and more.


\textbf{@hjanuschka} • \texttt{travel} \texttt{transport} \texttt{skill}

Real-time departures, disruptions, elevator status, and routing for Vienna's public transport.


\textbf{@George5562} • \texttt{automation} \texttt{browser} \texttt{parenting}

Automated UK school meal booking via ParentPay. Uses mouse coordinates for reliable table cell clicking.

\textbf{@julianengel} • \texttt{files} \texttt{r2} \texttt{presigned-urls}

Upload to Cloudflare R2/S3 and generate secure presigned download links. Perfect for remote OpenClaw instances.

\textbf{@coard} • \texttt{ios} \texttt{xcode} \texttt{testflight}

Built a complete iOS app with maps and voice recording, deployed to TestFlight entirely via Telegram chat.


\textbf{@AS} • \texttt{health} \texttt{oura} \texttt{calendar}

Personal AI health assistant integrating Oura ring data with calendar, appointments, and gym schedule.

\textbf{@adam91holt} • \texttt{multi-agent} \texttt{orchestration} \texttt{architecture} \texttt{manifesto}

14+ agents under one gateway with Opus 4.5 orchestrator delegating to Codex workers. Comprehensive \href{https://github.com/adam91holt/orchestrated-ai-articles}{technical write-up} covering the Dream Team roster, model selection, sandboxing, webhooks, heartbeats, and delegation flows. \href{https://github.com/adam91holt/clawdspace}{Clawdspace} for agent sandboxing. \href{https://adams-ai-journey.ghost.io/2026-the-year-of-the-orchestrator/}{Blog post}.

\textbf{@NessZerra} • \texttt{devtools} \texttt{linear} \texttt{cli} \texttt{issues}

CLI for Linear that integrates with agentic workflows (Claude Code, OpenClaw). Manage issues, projects, and workflows from the terminal. First external PR merged!

\textbf{@jules} • \texttt{messaging} \texttt{beeper} \texttt{cli} \texttt{automation}

Read, send, and archive messages via Beeper Desktop. Uses Beeper local MCP API so agents can manage all your chats (iMessage, WhatsApp, etc.) in one place.


\subsection{Automation \& Workflows}



\textbf{@antonplex} • \texttt{automation} \texttt{hardware} \texttt{air-quality}

Claude Code discovered and confirmed the purifier controls, then OpenClaw takes over to manage room air quality.


\textbf{@signalgaining} • \texttt{automation} \texttt{camera} \texttt{skill} \texttt{images}

Triggered by a roof camera: ask OpenClaw to snap a sky photo whenever it looks pretty — it designed a skill and took the shot.


\textbf{@buddyhadry} • \texttt{automation} \texttt{briefing} \texttt{images} \texttt{telegram}

A scheduled prompt generates a single "scene" image each morning (weather, tasks, date, favorite post/quote) via a OpenClaw persona.

\textbf{@joshp123} • \texttt{automation} \texttt{booking} \texttt{cli}

Playtomic availability checker + booking CLI. Never miss an open court again.


\textbf{Community} • \texttt{automation} \texttt{email} \texttt{pdf}

Collects PDFs from email, preps documents for tax consultant. Monthly accounting on autopilot.

\textbf{@davekiss} • \texttt{telegram} \texttt{website} \texttt{migration} \texttt{astro}

Rebuilt entire personal site via Telegram while watching Netflix — Notion → Astro, 18 posts migrated, DNS to Cloudflare. Never opened a laptop.

\textbf{@attol8} • \texttt{automation} \texttt{api} \texttt{skill}

Searches job listings, matches against CV keywords, and returns relevant opportunities with links. Built in 30 minutes using JSearch API.

\textbf{@jdrhyne} • \texttt{automation} \texttt{jira} \texttt{skill} \texttt{devtools}

OpenClaw connected to Jira, then generated a new skill on the fly (before it existed on ClawHub).

\textbf{@iamsubhrajyoti} • \texttt{automation} \texttt{todoist} \texttt{skill} \texttt{telegram}

Automated Todoist tasks and had OpenClaw generate the skill directly in Telegram chat.

\textbf{@bheem1798} • \texttt{finance} \texttt{browser} \texttt{automation}

Logs into TradingView via browser automation, screenshots charts, and performs technical analysis on demand. No API needed—just browser control.

\textbf{@henrymascot} • \texttt{slack} \texttt{automation} \texttt{support}

Watches company Slack channel, responds helpfully, and forwards notifications to Telegram. Autonomously fixed a production bug in a deployed app without being asked.


\subsection{Knowledge \& Memory}



\textbf{@joshp123} • \texttt{learning} \texttt{voice} \texttt{skill}

Chinese learning engine with pronunciation feedback and study flows via OpenClaw.


\textbf{Community} • \texttt{memory} \texttt{transcription} \texttt{indexing}

Ingests full WhatsApp exports, transcribes 1k+ voice notes, cross-checks with git logs, outputs linked markdown reports.

\textbf{@jamesbrooksco} • \texttt{search} \texttt{vector} \texttt{bookmarks}

Adds vector search to Karakeep bookmarks using Qdrant + OpenAI/Ollama embeddings.

\textbf{Community} • \texttt{memory} \texttt{beliefs} \texttt{self-model}

Separate memory manager that turns session files into memories → beliefs → evolving self model.


\subsection{Voice \& Phone}



\textbf{@alejandroOPI} • \texttt{voice} \texttt{vapi} \texttt{bridge}

Vapi voice assistant ↔ OpenClaw HTTP bridge. Near real-time phone calls with your agent.

\textbf{@obviyus} • \texttt{transcription} \texttt{multilingual} \texttt{skill}

Multi-lingual audio transcription via OpenRouter (Gemini, etc). Available on ClawHub.


\subsection{Infrastructure \& Deployment}



\textbf{@ngutman} • \texttt{homeassistant} \texttt{docker} \texttt{raspberry-pi}

OpenClaw gateway running on Home Assistant OS with SSH tunnel support and persistent state.

\textbf{ClawHub} • \texttt{homeassistant} \texttt{skill} \texttt{automation}

Control and automate Home Assistant devices via natural language.

\textbf{@openclaw} • \texttt{nix} \texttt{packaging} \texttt{deployment}

Batteries-included nixified OpenClaw configuration for reproducible deployments.

\textbf{ClawHub} • \texttt{calendar} \texttt{caldav} \texttt{skill}

Calendar skill using khal/vdirsyncer. Self-hosted calendar integration.


\subsection{Home \& Hardware}



\textbf{@joshp123} • \texttt{home} \texttt{nix} \texttt{grafana}

Nix-native home automation with OpenClaw as the interface, plus beautiful Grafana dashboards.


\textbf{@joshp123} • \texttt{vacuum} \texttt{iot} \texttt{plugin}

Control your Roborock robot vacuum through natural conversation.



\subsection{Community Projects}



\textbf{Community} • \texttt{marketplace} \texttt{astronomy} \texttt{webapp}

Full astronomy gear marketplace. Built with/around the OpenClaw ecosystem.


\hrulefill


\subsection{Submit Your Project}


Have something to share? We'd love to feature it!

Post in \href{https://discord.gg/clawd}{\#showcase on Discord} or \href{https://x.com/openclaw}{tweet @openclaw}
Tell us what it does, link to the repo/demo, share a screenshot if you have one
We'll add standout projects to this page


\clearpage

% === Source file: start/wizard-cli-automation.md ===
\section{CLI Automation}


Use \texttt{--non-interactive} to automate \texttt{openclaw onboard}.

\texttt{--json} does not imply non-interactive mode. Use \texttt{--non-interactive} (and \texttt{--workspace}) for scripts.

\subsection{Baseline non-interactive example}


\begin{lstlisting}[language=bash]
openclaw onboard --non-interactive \
  --mode local \
  --auth-choice apiKey \
  --anthropic-api-key "$ANTHROPIC_API_KEY" \
  --secret-input-mode plaintext \
  --gateway-port 18789 \
  --gateway-bind loopback \
  --install-daemon \
  --daemon-runtime node \
  --skip-skills
\end{lstlisting}


Add \texttt{--json} for a machine-readable summary.

Use \texttt{--secret-input-mode ref} to store env-backed refs in auth profiles instead of plaintext values.
Interactive selection between env refs and configured provider refs (\texttt{file} or \texttt{exec}) is available in the onboarding wizard flow.

In non-interactive \texttt{ref} mode, provider env vars must be set in the process environment.
Passing inline key flags without the matching env var now fails fast.

Example:

\begin{lstlisting}[language=bash]
openclaw onboard --non-interactive \
  --mode local \
  --auth-choice openai-api-key \
  --secret-input-mode ref \
  --accept-risk
\end{lstlisting}


\subsection{Provider-specific examples}


\begin{lstlisting}[language=bash]
    openclaw onboard --non-interactive \
      --mode local \
      --auth-choice gemini-api-key \
      --gemini-api-key "$GEMINI_API_KEY" \
      --gateway-port 18789 \
      --gateway-bind loopback
\end{lstlisting}

\begin{lstlisting}[language=bash]
    openclaw onboard --non-interactive \
      --mode local \
      --auth-choice zai-api-key \
      --zai-api-key "$ZAI_API_KEY" \
      --gateway-port 18789 \
      --gateway-bind loopback
\end{lstlisting}

\begin{lstlisting}[language=bash]
    openclaw onboard --non-interactive \
      --mode local \
      --auth-choice ai-gateway-api-key \
      --ai-gateway-api-key "$AI_GATEWAY_API_KEY" \
      --gateway-port 18789 \
      --gateway-bind loopback
\end{lstlisting}

\begin{lstlisting}[language=bash]
    openclaw onboard --non-interactive \
      --mode local \
      --auth-choice cloudflare-ai-gateway-api-key \
      --cloudflare-ai-gateway-account-id "your-account-id" \
      --cloudflare-ai-gateway-gateway-id "your-gateway-id" \
      --cloudflare-ai-gateway-api-key "$CLOUDFLARE_AI_GATEWAY_API_KEY" \
      --gateway-port 18789 \
      --gateway-bind loopback
\end{lstlisting}

\begin{lstlisting}[language=bash]
    openclaw onboard --non-interactive \
      --mode local \
      --auth-choice moonshot-api-key \
      --moonshot-api-key "$MOONSHOT_API_KEY" \
      --gateway-port 18789 \
      --gateway-bind loopback
\end{lstlisting}

\begin{lstlisting}[language=bash]
    openclaw onboard --non-interactive \
      --mode local \
      --auth-choice mistral-api-key \
      --mistral-api-key "$MISTRAL_API_KEY" \
      --gateway-port 18789 \
      --gateway-bind loopback
\end{lstlisting}

\begin{lstlisting}[language=bash]
    openclaw onboard --non-interactive \
      --mode local \
      --auth-choice synthetic-api-key \
      --synthetic-api-key "$SYNTHETIC_API_KEY" \
      --gateway-port 18789 \
      --gateway-bind loopback
\end{lstlisting}

\begin{lstlisting}[language=bash]
    openclaw onboard --non-interactive \
      --mode local \
      --auth-choice opencode-zen \
      --opencode-zen-api-key "$OPENCODE_API_KEY" \
      --gateway-port 18789 \
      --gateway-bind loopback
\end{lstlisting}

\begin{lstlisting}[language=bash]
    openclaw onboard --non-interactive \
      --mode local \
      --auth-choice custom-api-key \
      --custom-base-url "https://llm.example.com/v1" \
      --custom-model-id "foo-large" \
      --custom-api-key "$CUSTOM_API_KEY" \
      --custom-provider-id "my-custom" \
      --custom-compatibility anthropic \
      --gateway-port 18789 \
      --gateway-bind loopback
\end{lstlisting}


\texttt{--custom-api-key} is optional. If omitted, onboarding checks \texttt{CUSTOM\_API\_KEY}.

Ref-mode variant:

\begin{lstlisting}[language=bash]
    export CUSTOM_API_KEY="your-key"
    openclaw onboard --non-interactive \
      --mode local \
      --auth-choice custom-api-key \
      --custom-base-url "https://llm.example.com/v1" \
      --custom-model-id "foo-large" \
      --secret-input-mode ref \
      --custom-provider-id "my-custom" \
      --custom-compatibility anthropic \
      --gateway-port 18789 \
      --gateway-bind loopback
\end{lstlisting}


In this mode, onboarding stores \texttt{apiKey} as \texttt{\{ source: "env", provider: "default", id: "CUSTOM\_API\_KEY" \}}.


\subsection{Add another agent}


Use \texttt{openclaw agents add } to create a separate agent with its own workspace,
sessions, and auth profiles. Running without \texttt{--workspace} launches the wizard.

\begin{lstlisting}[language=bash]
openclaw agents add work \
  --workspace ~/.openclaw/workspace-work \
  --model openai/gpt-5.2 \
  --bind whatsapp:biz \
  --non-interactive \
  --json
\end{lstlisting}


What it sets:

\begin{itemize}
  \item \texttt{agents.list[].name}
  \item \texttt{agents.list[].workspace}
  \item \texttt{agents.list[].agentDir}
\end{itemize}


Notes:

\begin{itemize}
  \item Default workspaces follow \texttt{\textasciitilde{}/.openclaw/workspace-}.
  \item Add \texttt{bindings} to route inbound messages (the wizard can do this).
  \item Non-interactive flags: \texttt{--model}, \texttt{--agent-dir}, \texttt{--bind}, \texttt{--non-interactive}.
\end{itemize}


\subsection{Related docs}


\begin{itemize}
  \item Onboarding hub: \href{/start/wizard}{Onboarding Wizard (CLI)}
  \item Full reference: \href{/start/wizard-cli-reference}{CLI Onboarding Reference}
  \item Command reference: \href{/cli/onboard}{\texttt{openclaw onboard}}
\end{itemize}



\clearpage

% === Source file: start/wizard-cli-reference.md ===
\section{CLI Onboarding Reference}


This page is the full reference for \texttt{openclaw onboard}.
For the short guide, see \href{/start/wizard}{Onboarding Wizard (CLI)}.

\subsection{What the wizard does}


Local mode (default) walks you through:

\begin{itemize}
  \item Model and auth setup (OpenAI Code subscription OAuth, Anthropic API key or setup token, plus MiniMax, GLM, Moonshot, and AI Gateway options)
  \item Workspace location and bootstrap files
  \item Gateway settings (port, bind, auth, tailscale)
  \item Channels and providers (Telegram, WhatsApp, Discord, Google Chat, Mattermost plugin, Signal)
  \item Daemon install (LaunchAgent or systemd user unit)
  \item Health check
  \item Skills setup
\end{itemize}


Remote mode configures this machine to connect to a gateway elsewhere.
It does not install or modify anything on the remote host.

\subsection{Local flow details}


\begin{itemize}
  \item If \texttt{\textasciitilde{}/.openclaw/openclaw.json} exists, choose Keep, Modify, or Reset.
  \item Re-running the wizard does not wipe anything unless you explicitly choose Reset (or pass \texttt{--reset}).
  \item CLI \texttt{--reset} defaults to \texttt{config+creds+sessions}; use \texttt{--reset-scope full} to also remove workspace.
  \item If config is invalid or contains legacy keys, the wizard stops and asks you to run \texttt{openclaw doctor} before continuing.
  \item Reset uses \texttt{trash} and offers scopes:
  \item Config only
  \item Config + credentials + sessions
  \item Full reset (also removes workspace)
\end{itemize}

\begin{itemize}
  \item Full option matrix is in \href{\#auth-and-model-options}{Auth and model options}.
\end{itemize}

\begin{itemize}
  \item Default \texttt{\textasciitilde{}/.openclaw/workspace} (configurable).
  \item Seeds workspace files needed for first-run bootstrap ritual.
  \item Workspace layout: \href{/concepts/agent-workspace}{Agent workspace}.
\end{itemize}

\begin{itemize}
  \item Prompts for port, bind, auth mode, and tailscale exposure.
  \item Recommended: keep token auth enabled even for loopback so local WS clients must authenticate.
  \item Disable auth only if you fully trust every local process.
  \item Non-loopback binds still require auth.
\end{itemize}

\begin{itemize}
  \item \href{/channels/whatsapp}{WhatsApp}: optional QR login
  \item \href{/channels/telegram}{Telegram}: bot token
  \item \href{/channels/discord}{Discord}: bot token
  \item \href{/channels/googlechat}{Google Chat}: service account JSON + webhook audience
  \item \href{/channels/mattermost}{Mattermost} plugin: bot token + base URL
  \item \href{/channels/signal}{Signal}: optional \texttt{signal-cli} install + account config
  \item \href{/channels/bluebubbles}{BlueBubbles}: recommended for iMessage; server URL + password + webhook
  \item \href{/channels/imessage}{iMessage}: legacy \texttt{imsg} CLI path + DB access
  \item DM security: default is pairing. First DM sends a code; approve via
\end{itemize}

\texttt{openclaw pairing approve  } or use allowlists.
\begin{itemize}
  \item macOS: LaunchAgent
  \item Requires logged-in user session; for headless, use a custom LaunchDaemon (not shipped).
  \item Linux and Windows via WSL2: systemd user unit
  \item Wizard attempts \texttt{loginctl enable-linger } so gateway stays up after logout.
  \item May prompt for sudo (writes \texttt{/var/lib/systemd/linger}); it tries without sudo first.
  \item Runtime selection: Node (recommended; required for WhatsApp and Telegram). Bun is not recommended.
\end{itemize}

\begin{itemize}
  \item Starts gateway (if needed) and runs \texttt{openclaw health}.
  \item \texttt{openclaw status --deep} adds gateway health probes to status output.
\end{itemize}

\begin{itemize}
  \item Reads available skills and checks requirements.
  \item Lets you choose node manager: npm or pnpm (bun not recommended).
  \item Installs optional dependencies (some use Homebrew on macOS).
\end{itemize}

\begin{itemize}
  \item Summary and next steps, including iOS, Android, and macOS app options.
\end{itemize}


If no GUI is detected, the wizard prints SSH port-forward instructions for the Control UI instead of opening a browser.
If Control UI assets are missing, the wizard attempts to build them; fallback is \texttt{pnpm ui:build} (auto-installs UI deps).

\subsection{Remote mode details}


Remote mode configures this machine to connect to a gateway elsewhere.

Remote mode does not install or modify anything on the remote host.

What you set:

\begin{itemize}
  \item Remote gateway URL (\texttt{ws://...})
  \item Token if remote gateway auth is required (recommended)
\end{itemize}


\begin{itemize}
  \item If gateway is loopback-only, use SSH tunneling or a tailnet.
  \item Discovery hints:
  \item macOS: Bonjour (\texttt{dns-sd})
  \item Linux: Avahi (\texttt{avahi-browse})
\end{itemize}


\subsection{Auth and model options}


Uses \texttt{ANTHROPIC\_API\_KEY} if present or prompts for a key, then saves it for daemon use.
\begin{itemize}
  \item macOS: checks Keychain item "Claude Code-credentials"
  \item Linux and Windows: reuses \texttt{\textasciitilde{}/.claude/.credentials.json} if present
\end{itemize}


On macOS, choose "Always Allow" so launchd starts do not block.

Run \texttt{claude setup-token} on any machine, then paste the token.
You can name it; blank uses default.
If \texttt{\textasciitilde{}/.codex/auth.json} exists, the wizard can reuse it.
Browser flow; paste \texttt{code\#state}.

Sets \texttt{agents.defaults.model} to \texttt{openai-codex/gpt-5.3-codex} when model is unset or \texttt{openai/*}.

Uses \texttt{OPENAI\_API\_KEY} if present or prompts for a key, then stores the credential in auth profiles.

Sets \texttt{agents.defaults.model} to \texttt{openai/gpt-5.1-codex} when model is unset, \texttt{openai/\textit{}, or \texttt{openai-codex/}}.

Prompts for \texttt{XAI\_API\_KEY} and configures xAI as a model provider.
Prompts for \texttt{OPENCODE\_API\_KEY} (or \texttt{OPENCODE\_ZEN\_API\_KEY}).
Setup URL: \href{https://opencode.ai/auth}{opencode.ai/auth}.
Stores the key for you.
Prompts for \texttt{AI\_GATEWAY\_API\_KEY}.
More detail: \href{/providers/vercel-ai-gateway}{Vercel AI Gateway}.
Prompts for account ID, gateway ID, and \texttt{CLOUDFLARE\_AI\_GATEWAY\_API\_KEY}.
More detail: \href{/providers/cloudflare-ai-gateway}{Cloudflare AI Gateway}.
Config is auto-written.
More detail: \href{/providers/minimax}{MiniMax}.
Prompts for \texttt{SYNTHETIC\_API\_KEY}.
More detail: \href{/providers/synthetic}{Synthetic}.
Moonshot (Kimi K2) and Kimi Coding configs are auto-written.
More detail: \href{/providers/moonshot}{Moonshot AI (Kimi + Kimi Coding)}.
Works with OpenAI-compatible and Anthropic-compatible endpoints.

Interactive onboarding supports the same API key storage choices as other provider API key flows:
\begin{itemize}
  \item \textbf{Paste API key now} (plaintext)
  \item \textbf{Use secret reference} (env ref or configured provider ref, with preflight validation)
\end{itemize}


Non-interactive flags:
\begin{itemize}
  \item \texttt{--auth-choice custom-api-key}
  \item \texttt{--custom-base-url}
  \item \texttt{--custom-model-id}
  \item \texttt{--custom-api-key} (optional; falls back to \texttt{CUSTOM\_API\_KEY})
  \item \texttt{--custom-provider-id} (optional)
  \item \texttt{--custom-compatibility } (optional; default \texttt{openai})
\end{itemize}


Leaves auth unconfigured.

Model behavior:

\begin{itemize}
  \item Pick default model from detected options, or enter provider and model manually.
  \item Wizard runs a model check and warns if the configured model is unknown or missing auth.
\end{itemize}


Credential and profile paths:

\begin{itemize}
  \item OAuth credentials: \texttt{\textasciitilde{}/.openclaw/credentials/oauth.json}
  \item Auth profiles (API keys + OAuth): \texttt{\textasciitilde{}/.openclaw/agents//agent/auth-profiles.json}
\end{itemize}


API key storage mode:

\begin{itemize}
  \item Default onboarding behavior persists API keys as plaintext values in auth profiles.
  \item \texttt{--secret-input-mode ref} enables reference mode instead of plaintext key storage.
\end{itemize}

In interactive onboarding, you can choose either:
\begin{itemize}
  \item environment variable ref (for example \texttt{keyRef: \{ source: "env", provider: "default", id: "OPENAI\_API\_KEY" \}})
  \item configured provider ref (\texttt{file} or \texttt{exec}) with provider alias + id
  \item Interactive reference mode runs a fast preflight validation before saving.
  \item Env refs: validates variable name + non-empty value in the current onboarding environment.
  \item Provider refs: validates provider config and resolves the requested id.
  \item If preflight fails, onboarding shows the error and lets you retry.
  \item In non-interactive mode, \texttt{--secret-input-mode ref} is env-backed only.
  \item Set the provider env var in the onboarding process environment.
  \item Inline key flags (for example \texttt{--openai-api-key}) require that env var to be set; otherwise onboarding fails fast.
  \item For custom providers, non-interactive \texttt{ref} mode stores \texttt{models.providers..apiKey} as \texttt{\{ source: "env", provider: "default", id: "CUSTOM\_API\_KEY" \}}.
  \item In that custom-provider case, \texttt{--custom-api-key} requires \texttt{CUSTOM\_API\_KEY} to be set; otherwise onboarding fails fast.
  \item Existing plaintext setups continue to work unchanged.
\end{itemize}


Headless and server tip: complete OAuth on a machine with a browser, then copy
\texttt{\textasciitilde{}/.openclaw/credentials/oauth.json} (or \texttt{\$OPENCLAW\_STATE\_DIR/credentials/oauth.json})
to the gateway host.

\subsection{Outputs and internals}


Typical fields in \texttt{\textasciitilde{}/.openclaw/openclaw.json}:

\begin{itemize}
  \item \texttt{agents.defaults.workspace}
  \item \texttt{agents.defaults.model} / \texttt{models.providers} (if Minimax chosen)
  \item \texttt{tools.profile} (local onboarding defaults to \texttt{"messaging"} when unset; existing explicit values are preserved)
  \item \texttt{gateway.*} (mode, bind, auth, tailscale)
  \item \texttt{session.dmScope} (local onboarding defaults this to \texttt{per-channel-peer} when unset; existing explicit values are preserved)
  \item \texttt{channels.telegram.botToken}, \texttt{channels.discord.token}, \texttt{channels.signal.\textit{}, \texttt{channels.imessage.}}
  \item Channel allowlists (Slack, Discord, Matrix, Microsoft Teams) when you opt in during prompts (names resolve to IDs when possible)
  \item \texttt{skills.install.nodeManager}
  \item \texttt{wizard.lastRunAt}
  \item \texttt{wizard.lastRunVersion}
  \item \texttt{wizard.lastRunCommit}
  \item \texttt{wizard.lastRunCommand}
  \item \texttt{wizard.lastRunMode}
\end{itemize}


\texttt{openclaw agents add} writes \texttt{agents.list[]} and optional \texttt{bindings}.

WhatsApp credentials go under \texttt{\textasciitilde{}/.openclaw/credentials/whatsapp//}.
Sessions are stored under \texttt{\textasciitilde{}/.openclaw/agents//sessions/}.

Some channels are delivered as plugins. When selected during onboarding, the wizard
prompts to install the plugin (npm or local path) before channel configuration.

Gateway wizard RPC:

\begin{itemize}
  \item \texttt{wizard.start}
  \item \texttt{wizard.next}
  \item \texttt{wizard.cancel}
  \item \texttt{wizard.status}
\end{itemize}


Clients (macOS app and Control UI) can render steps without re-implementing onboarding logic.

Signal setup behavior:

\begin{itemize}
  \item Downloads the appropriate release asset
  \item Stores it under \texttt{\textasciitilde{}/.openclaw/tools/signal-cli//}
  \item Writes \texttt{channels.signal.cliPath} in config
  \item JVM builds require Java 21
  \item Native builds are used when available
  \item Windows uses WSL2 and follows Linux signal-cli flow inside WSL
\end{itemize}


\subsection{Related docs}


\begin{itemize}
  \item Onboarding hub: \href{/start/wizard}{Onboarding Wizard (CLI)}
  \item Automation and scripts: \href{/start/wizard-cli-automation}{CLI Automation}
  \item Command reference: \href{/cli/onboard}{\texttt{openclaw onboard}}
\end{itemize}



\clearpage

% === Source file: start/wizard.md ===
\section{Onboarding Wizard (CLI)}


The onboarding wizard is the \textbf{recommended} way to set up OpenClaw on macOS,
Linux, or Windows (via WSL2; strongly recommended).
It configures a local Gateway or a remote Gateway connection, plus channels, skills,
and workspace defaults in one guided flow.

\begin{lstlisting}[language=bash]
openclaw onboard
\end{lstlisting}


Fastest first chat: open the Control UI (no channel setup needed). Run
\texttt{openclaw dashboard} and chat in the browser. Docs: \href{/web/dashboard}{Dashboard}.

To reconfigure later:

\begin{lstlisting}[language=bash]
openclaw configure
openclaw agents add <name>
\end{lstlisting}


\texttt{--json} does not imply non-interactive mode. For scripts, use \texttt{--non-interactive}.

Recommended: set up a Brave Search API key so the agent can use \texttt{web\_search}
(\texttt{web\_fetch} works without a key). Easiest path: \texttt{openclaw configure --section web}
which stores \texttt{tools.web.search.apiKey}. Docs: \href{/tools/web}{Web tools}.

\subsection{QuickStart vs Advanced}


The wizard starts with \textbf{QuickStart} (defaults) vs \textbf{Advanced} (full control).

\begin{itemize}
  \item Local gateway (loopback)
  \item Workspace default (or existing workspace)
  \item Gateway port \textbf{18789}
  \item Gateway auth \textbf{Token} (auto‑generated, even on loopback)
  \item Tool policy default for new local setups: \texttt{tools.profile: "messaging"} (existing explicit profile is preserved)
  \item DM isolation default: local onboarding writes \texttt{session.dmScope: "per-channel-peer"} when unset. Details: \href{/start/wizard-cli-reference\#outputs-and-internals}{CLI Onboarding Reference}
  \item Tailscale exposure \textbf{Off}
  \item Telegram + WhatsApp DMs default to \textbf{allowlist} (you'll be prompted for your phone number)
\end{itemize}

\begin{itemize}
  \item Exposes every step (mode, workspace, gateway, channels, daemon, skills).
\end{itemize}


\subsection{What the wizard configures}


\textbf{Local mode (default)} walks you through these steps:

\begin{enumerate}
  \item \textbf{Model/Auth} — choose any supported provider/auth flow (API key, OAuth, or setup-token), including Custom Provider
\end{enumerate}

(OpenAI-compatible, Anthropic-compatible, or Unknown auto-detect). Pick a default model.
Security note: if this agent will run tools or process webhook/hooks content, prefer the strongest latest-generation model available and keep tool policy strict. Weaker/older tiers are easier to prompt-inject.
For non-interactive runs, \texttt{--secret-input-mode ref} stores env-backed refs in auth profiles instead of plaintext API key values.
In non-interactive \texttt{ref} mode, the provider env var must be set; passing inline key flags without that env var fails fast.
In interactive runs, choosing secret reference mode lets you point at either an environment variable or a configured provider ref (\texttt{file} or \texttt{exec}), with a fast preflight validation before saving.
\begin{enumerate}
  \item \textbf{Workspace} — Location for agent files (default \texttt{\textasciitilde{}/.openclaw/workspace}). Seeds bootstrap files.
  \item \textbf{Gateway} — Port, bind address, auth mode, Tailscale exposure.
  \item \textbf{Channels} — WhatsApp, Telegram, Discord, Google Chat, Mattermost, Signal, BlueBubbles, or iMessage.
  \item \textbf{Daemon} — Installs a LaunchAgent (macOS) or systemd user unit (Linux/WSL2).
  \item \textbf{Health check} — Starts the Gateway and verifies it's running.
  \item \textbf{Skills} — Installs recommended skills and optional dependencies.
\end{enumerate}


Re-running the wizard does \textbf{not} wipe anything unless you explicitly choose \textbf{Reset} (or pass \texttt{--reset}).
CLI \texttt{--reset} defaults to config, credentials, and sessions; use \texttt{--reset-scope full} to include workspace.
If the config is invalid or contains legacy keys, the wizard asks you to run \texttt{openclaw doctor} first.

\textbf{Remote mode} only configures the local client to connect to a Gateway elsewhere.
It does \textbf{not} install or change anything on the remote host.

\subsection{Add another agent}


Use \texttt{openclaw agents add } to create a separate agent with its own workspace,
sessions, and auth profiles. Running without \texttt{--workspace} launches the wizard.

What it sets:

\begin{itemize}
  \item \texttt{agents.list[].name}
  \item \texttt{agents.list[].workspace}
  \item \texttt{agents.list[].agentDir}
\end{itemize}


Notes:

\begin{itemize}
  \item Default workspaces follow \texttt{\textasciitilde{}/.openclaw/workspace-}.
  \item Add \texttt{bindings} to route inbound messages (the wizard can do this).
  \item Non-interactive flags: \texttt{--model}, \texttt{--agent-dir}, \texttt{--bind}, \texttt{--non-interactive}.
\end{itemize}


\subsection{Full reference}


For detailed step-by-step breakdowns, non-interactive scripting, Signal setup,
RPC API, and a full list of config fields the wizard writes, see the
\href{/reference/wizard}{Wizard Reference}.

\subsection{Related docs}


\begin{itemize}
  \item CLI command reference: \href{/cli/onboard}{\texttt{openclaw onboard}}
  \item Onboarding overview: \href{/start/onboarding-overview}{Onboarding Overview}
  \item macOS app onboarding: \href{/start/onboarding}{Onboarding}
  \item Agent first-run ritual: \href{/start/bootstrapping}{Agent Bootstrapping}
\end{itemize}



\clearpage
